\label{sec:discussion}

\subsection{Interpretation of Findings}

The central finding of this paper---that enacted maps achieve pairing rates far below
the random baseline while \textsc{bisect} maps fall within the baseline---has a
straightforward interpretation: incumbent mapmakers use incumbency data to draw district
lines that guarantee each sitting member a home district. The algorithm cannot do this
because it does not receive incumbency data. The resulting difference in pairing rates
is not attributable to any difference in geographic constraints (both maps must satisfy
the same contiguity, population, and compactness requirements); it is attributable
entirely to the enacted mapmakers' access to and use of incumbent address data.

This framing has an important implication: enacted maps' low pairing rates are not a
gift from geography. They are a choice. A mapmaker with no incumbency data---or with
an explicit instruction not to use incumbency data---would produce pairing rates in the
range of the algorithmic baseline, not the enacted zero. The incumbency protection
premium quantifies the cost of this choice: to achieve zero pairings in NC, mapmakers
must draw district boundaries that deviate from the geographic optimum in ways that
benefit specific incumbents, at the cost of whatever other criteria would have been
served by the geographically optimal boundary.

\subsection{Same-Party vs.\ Cross-Party Pairing}

The pairing events under \textsc{bisect} maps in this study are predominantly
same-party: both NC pairings in the single-seed result involve Republican incumbents
whose home tracts are geographically proximate. This is not surprising: geographic
proximity of home tracts is correlated with same-party membership in states with
strong geographic partisan sorting, because incumbents of the same party tend to
represent geographically adjacent districts and therefore tend to live in adjacent
or nearby tracts.

Cross-party pairing---where a Republican and a Democratic incumbent find themselves
in the same new district---is more likely to occur in states with significant partisan
mixing at the district boundary, which is more common in competitive states like WI.
In our sample, the single WI pairing event under \textsc{bisect}$^\dagger$ involves
incumbents from different parties.

For litigation purposes, same-party pairing and cross-party pairing have different
political valences. Same-party pairing produces a primary contest with uncertain
implications for the general election. Cross-party pairing produces an automatic general
election contest between sitting incumbents, with significant media attention and
fundraising implications. Courts have not distinguished between these types in
evaluating whether pairing constitutes an unconstitutional burden on incumbents; neither
type is constitutionally prohibited.

\subsection{Legal Framework: Incumbency as a Permissible Factor}

The legal framework governing incumbent pairing in redistricting is established by
three key propositions. First, incumbent pairing is not constitutionally prohibited.
No constitutional provision prevents a redistricting plan from placing two incumbents
in the same district; the First Amendment and the Equal Protection Clause protect
voters, not the electoral interests of sitting legislators~\citep{rucho2019}.

Second, incumbent protection is permissible but not required. As the Supreme Court
noted in \textit{Karcher v.\ Daggett}\footnote{\textit{Karcher v.\ Daggett}, 462
U.S.\ 725, 740 (1983).} and \textit{Cox v.\ Larios}\footnote{\textit{Cox v.\ Larios},
542 U.S.\ 947 (2004).}, a mapmaker may consider incumbency as one factor among
many, but may not justify unconstitutional population deviations or other constitutional
violations on the sole basis of incumbency protection.

Third, an algorithm that ignores incumbency is constitutionally sound. There is no
federal constitutional right to a redistricting process that considers incumbency;
voters have no cognisable injury from a plan that produces pairing rates consistent
with random redistricting. State statutory requirements to avoid ``undue'' incumbent
disruption may impose an additional obligation, but these requirements have generally
been interpreted as requiring good-faith consideration of incumbency, not as mandating
any specific pairing rate.

The practical implication for expert witnesses is clear: a \textsc{bisect} map that
produces pairing rates consistent with the random baseline is constitutionally sound.
The fact that the enacted map achieves a lower pairing rate does not mean that the
algorithmic map is constitutionally defective; it means that the enacted map chose
to use incumbency data to reduce pairing, a choice the algorithm did not make.

\subsection{Evidence of Strategic Mapmaking}

The incumbency protection premium provides a useful quantitative tool for redistricting
litigation. A large positive IPP---the enacted map produces far fewer pairings than
geography would imply---is circumstantial evidence that mapmakers used incumbent
location data in drawing district lines. This evidence is strongest when combined with
evidence that other traditional redistricting criteria (compactness, community
preservation) were subordinated to achieve the low pairing rate~\citep{mcghee2014measuring}.

For NC and WI, where the enacted IPP is 0.175 or greater, the evidence is strong:
a map that achieves zero pairings among 8 or 13 incumbents in a state where random
redistricting would produce multiple pairings cannot plausibly have been drawn without
reference to incumbent home addresses. For TX, where the enacted IPP is 0.127, the
evidence is somewhat weaker---a low-pairing outcome is less improbable in a large
delegation---but still below the random baseline under the preliminary single-seed binomial test (exact permutation tests are a Phase~2 extension).
